\documentclass[12pt]{article}
\usepackage{anysize}
\usepackage{xltxtra}
\usepackage{xunicode}
\usepackage{graphicx}
\usepackage{color}
\usepackage{tikz}
\usepackage{minted}
\usepackage{listings}
\usepackage{enumitem}
\usepackage{framed}
\usepackage{relsize}
\usepackage{float}
\usepackage{algpseudocode}
\usepackage[hyphens]{url}

\usepackage{color,hyperref}

\usepackage{fancyhdr}
\usepackage{lastpage}
\usepackage{fontspec}

\linespread{1.5}

\definecolor{darkblue}{rgb}{0.0,0.0,0.3}
\hypersetup{colorlinks,breaklinks,
        linkcolor=darkblue,urlcolor=darkblue,
            anchorcolor=black,citecolor=black}

\def\name{Grigorios Lyras}
\def\IF{\textbf{if}}
\def\THEN{\textbf{then}}
\def\SAYS{\textbf{says}}
\def\SPEAKSFOR{$\Rightarrow$}
\def\FOR{\textbf{for}}
\def\AND{\textbf{and}}

\setlength{\headheight}{15.2pt}
\pagestyle{fancy}{ %
    \fancyhf{} % remove everything
    \renewcommand{\headrulewidth}{0pt} % remove lines as well
    \renewcommand{\footrulewidth}{0.5pt}
    \lfoot{\name}
    \rfoot{Page \thepage\ of \pageref{LastPage}}

\renewcommand{\labelenumi}{(\alph{enumi})}
\defaultfontfeatures{Ligatures=TeX} % To support LaTeX quoting style
\setmainfont[Mapping=TeX-text]{CMU Serif}


\begin{document}

\begin{titlepage}

\begin{tikzpicture}[remember picture, overlay]
      \node [anchor=north east, inner sep=0pt]  at (current page.north east)
           {\includegraphics[height=3cm]{title/ucl_logo}};
       \end{tikzpicture}
\begin{center}
\begin{LARGE}\textbf{
University College London\\}\end{LARGE}
\vspace{2cm}
\begin{Large}
MSc Networked Computer Systems\\
Dronkeys experiment report\\
\end{Large}
\large{Academic Year 2013-2014}\\
\vspace{1cm}
\begin{tabular}{l r}
\Large{Grigorios Lyras}&
\large{SID: 13047325}\\
\end{tabular}\\
\vfill
\large\today\\
\end{center}
\end{titlepage}



\section{Introduction}

This report is to describe the experiment that took place at the campus of the
Royal Veterinary College on the 7th of August in collaboration with Andrew
Symington and Ian <XXX: fill name here>


\section{Experiment Description}
The purpose of the experiment was to gather data on the behaviour of the radio
channel in an Unmanned Air Vehicle environment. The data to be gathered were
radio measurements in terms of Receiver Signal Strength Indicator (RSSI) and Bit
Error Rate (BER). The data gathered were also measurements from the UAV sensors
for acceleration, speed and orientation along with magnetometer readings.
Finally these are to be fused with Global Position System data. The research
focus is to identify and model the correlation between magnetometer and relative
position and orientation data and radio measurements. The radio measurements
were taken on WiFi and ZigBee on non overlapping channels.

\section{Experiment Log}
During the day we used used two <XXX: insert drone model here> to conduct the
experiment. These had attached on them a Fit2PC Rev1.4 which was equipped with a
PCIe Wireless Card <XXX:insert model here>. The UAVs were also carrying two
prototyping boards which held an OrisenStamp Model 2013.

We arrived at RVC at 9:00 and started preparing the UAVs for their flight.
After mounting the boards on the UAVs, the GPS antenna at the top for best
reception and the ultra sonic module at the bottom for better accuracy we
proceeded to test the setup through the preflight checklist. Andrew had equipped
them with two Fit2PCs however we soon realised that one would not boot. At that
point we proceeded with a replacement of the Fit2PC with a working one. During
the swap of the Fit2PC half the team proceeded in the field to setup the base
station and start capturing GPS data, both RAW measurements and the solution
that rtknavi produced were captured. The capture begun at 11:19 am. In order to
verify and double check the capture settings the capture was stopped twice
however the settings were not altered therefore the data captured can simply be
appended and treated as one log.

When the UAVs were fitted with working hardware we brought them in the field and
started the final preflight checks to verify that they booted properly and they
were capturing data properly. The OrisenStamp was capturing properly however the
Fit2PCs were not booting properly again. We moved them back to the barn to
identify the bug and fix the issue. We then decided to give up our slot to the
next group so that we make better use of the day.

After some debugging in the bootloader configuration we fixed the issue and
proceeded in the field to perform our experiment. At <insert timestamp here> we
started our experiment again. Due to the fact that we were testing the UAVs
beforehand the batteries had a drop in their load. This allowed the UAVs to
fly for a duration of approximately 5 minutes. When Andrew and Ian landed
the quadcopters we realised that one of the batteries had pushed and unclipped
the external ZigBee antenna on the moving UAV (Golf). We unplugged the power
supply, put in new batteries, plugged in the antenna and restarted the
experiment from the top. At that point the SD cards were not erased, instead the
log kept being appended.

During the second attempt we flew the UAVs at a height of 13 meters. One was
hovering and slightly moving due to wind while the other was trying to cover a
box area under it as shown in Figure <XXX:insert figure here>. The flight time
was in the order of 10 minutes before the stationary UAV (Foxtrot) gave a
battery warning. Both Vehicles were landed and Foxtrot was shutdown. Our
experiment was complete. At that point Mahdi requested to take some extra
measurements using the other UAV (Golf). Golf was still powered up while we took
some measurements lifting the UAV to roughtly two meters and then slowly taking
it down by hand. Then we replaced <XXX: Johann please verify> Golf's battery
pack and Andrew made some low altitude hovering for Mahdi's experiment. After
that because the hovering was not stable enough due to wind conditions, we
unplugged Golf's power supply and moved back into the barn. There Golf was
powered up, (once to the best of my knowledge), and again Andrew flew it low for
a few minutes. With that our day was concluded. We passed the data to Andrew
who uploaded them on a dropbox folder found <XXX:here>.

\end{document}
